% !TeX root = ../Report.tex
\chapter{Future Work}
\label{sec:future work}

The results obtained in this evaluation provide several directions for future research. These directions concern both the evolution of the FUNES architecture and the broader evaluation methodology, including the underlying language models, benchmark design, and validation in operational settings. The main objective is to determine whether the improvements observed in the current benchmark generalize beyond the considered scenario and to identify which architectural mechanisms are responsible for the observed behaviour.

\section{Architectural Improvements}

A first direction concerns the extension of the retrieval and orchestration mechanisms towards more explicit requirement-aware reasoning. The current architecture focuses primarily on identifying and retrieving the information required to answer a query, while the evaluation results show that some failures occur even when the relevant evidence has already been retrieved. In particular, the results for Q19 indicate that the system may identify the relevant anomalies without fully mapping them to all the entities, attributes, and temporal information explicitly requested by the user. Future work could therefore introduce a requirement-tracking mechanism that decomposes the query into explicit answer requirements and verifies, during the synthesis and validation stages, whether each requirement has been addressed. Such a mechanism could verify, for example, that all requested activities, datetimes, affected resources, and operational relationships are explicitly represented in the final response.

A related direction concerns the development of more adaptive retrieval strategies. The current architecture relies on the Planner Agent to determine which information sources should be queried, while subsequent components process the retrieved evidence. Future versions of FUNES could instead adopt an iterative retrieval strategy in which the system evaluates whether the evidence collected so far is sufficient to answer the query and dynamically determines whether additional information should be retrieved. This approach could reduce both under-retrieval, where relevant evidence is missing, and over-retrieval, where excessive information increases the complexity of the subsequent reasoning process. Such a mechanism could be particularly relevant for complex operational investigations in which the relevance of additional evidence becomes apparent only after an initial retrieval step.

The representation of evidence could also be extended through explicit provenance and relationship tracking. Rather than providing the Synthesizer Agent with a collection of retrieved facts, future implementations could maintain an evidence graph representing relationships between anomalies, assets, operational activities, tickets, planning events, and observed impacts. Each element could retain its source and temporal information, allowing the final reasoning stage to distinguish directly observed facts from inferred relationships. This could improve both answer traceability and the ability of the Critic Agent to verify whether the conclusions generated by the system are supported by the available evidence.

The Critic Agent could consequently be extended beyond the current focus on evidence sufficiency and retrieval refinement. Future implementations could introduce requirement-level validation, factual consistency checks, temporal consistency checks, and provenance verification. For example, the Critic could explicitly verify that every activity identified as affected by an anomaly is supported by the available planning evidence and that the corresponding datetime is correctly reported. This would shift part of the validation process from checking whether sufficient information has been retrieved to checking whether the retrieved information has been correctly transformed into the requested operational conclusion.

Another possible extension concerns uncertainty and confidence representation. Operational decision-support systems must distinguish between confirmed observations, supported inferences, and hypotheses that require further investigation. FUNES could therefore associate confidence or evidence-status information with retrieved facts and generated conclusions. This would allow the Synthesizer Agent to explicitly communicate when a causal relationship is not established by the available evidence, reducing the risk of presenting plausible operational interpretations as confirmed facts.

Finally, future work could investigate adaptive and cost-aware orchestration. The current architecture prioritizes information relevance and answer quality, but operational deployment also requires controlling latency, computational cost, number of tool calls, and token consumption. The orchestration strategy could therefore be optimized jointly for answer quality and resource consumption. This would allow the system to determine when an additional retrieval or refinement cycle provides sufficient expected benefit to justify its computational cost.

\section{LLM and Reasoning Strategies}

A second direction concerns the interaction between the FUNES architecture and the underlying language model. The current evaluation compares FUNES with a single baseline based on Llama 3.3 70B Instruct. Although this provides a controlled comparison, it does not establish whether the observed benefits are independent of the characteristics of the underlying model. Future work should therefore evaluate FUNES with multiple LLM families and model sizes, including both smaller and larger models.

In particular, an interesting research direction would be to investigate whether explicit information orchestration can compensate, at least partially, for the lower reasoning capacity of smaller language models. Experiments could compare a large LLM operating on the complete operational context with a smaller LLM operating within the FUNES architecture. Such a comparison would provide evidence on whether the architectural mechanisms can reduce the dependence on increasingly large models for operational decision-support tasks.

Future evaluations could also investigate different reasoning and prompting strategies within the same FUNES architecture. This could include structured intermediate representations, explicit reasoning plans, constrained output formats, and different approaches to final answer synthesis. The objective would not simply be to identify the highest-scoring configuration, but to determine whether the architectural separation between information acquisition and reasoning remains beneficial across different LLM configurations.

Another important direction concerns the evaluation of the LLM-based judge itself. The current benchmark relies on an evaluator model to assign scores according to the defined criteria. Future work should validate the reliability of this evaluation methodology through comparison with human domain experts and, where possible, multiple independent evaluator models. Agreement between different evaluators could be analysed to identify criteria for which automated evaluation is stable and criteria for which human judgement remains particularly important.

The evaluation could also be extended from absolute criterion scoring to pairwise comparison. In addition to assigning independent scores to each system, evaluators could determine which of two responses is more correct or operationally useful with respect to the same ground truth. Combining absolute scoring and pairwise preference could provide a more robust evaluation of relative system performance and reduce the sensitivity of the results to the boundaries of the discrete scoring scale.

\section{Benchmark and Evaluation Extensions}

The current benchmark is intentionally limited to a controlled operational scenario and a relatively small number of questions. A primary direction for future work is therefore the construction of a larger and more diverse benchmark. Increasing the number of questions would reduce the influence of individual cases such as Q16 and Q19 and would allow performance differences to be analysed with greater statistical confidence.

The benchmark should also include a broader range of operational scenarios, ground stations, satellites, anomalies, planning constraints, and combinations of operational events. This would allow the evaluation to determine whether the observed behaviour is specific to the current scenario or persists across different configurations of the ground segment.

A more systematic definition of task complexity would also be beneficial. Instead of assigning questions to broad easy, medium, and hard categories, future benchmarks could define complexity along multiple dimensions, such as number of information sources, number of reasoning steps, number of entities involved, temporal dependencies, degree of ambiguity, and requirement for operational recommendations. This would allow performance to be analysed as a function of specific task characteristics rather than a single categorical notion of difficulty.

Future benchmark versions should also include a larger proportion of multi-source and multi-hop questions. These tasks would provide a more direct evaluation of the capabilities that motivate the FUNES architecture, particularly the ability to identify relevant sources, correlate heterogeneous evidence, and preserve relationships across multiple retrieval steps. The benchmark could explicitly control the number of sources required to answer each question, enabling the relationship between source heterogeneity and system performance to be investigated.

Another important extension would be the introduction of noisy and adversarial contexts. Irrelevant information, redundant records, temporally similar events, and operational anomalies that are correlated but not causally related could be deliberately introduced into the context. This would allow the robustness of FUNES and the baseline to be evaluated under controlled levels of information overload and ambiguity. Such experiments would provide a stronger test of the hypothesis that explicit information selection can reduce errors caused by irrelevant contextual information.

The benchmark should also include questions for which the available evidence is intentionally insufficient to determine a unique answer. These cases are particularly relevant for operational decision support, since a reliable system should be able to distinguish between a supported conclusion and a situation requiring additional investigation. Evaluating the ability to explicitly identify insufficient evidence would therefore complement the current focus on answer correctness.

Temporal reasoning represents another important benchmark dimension. Satellite ground-segment operations are inherently time-dependent, and future benchmark scenarios should explicitly test relationships between anomaly occurrence, planned activities, maintenance windows, acquisition opportunities, and subsequent operational impacts. Questions could require systems to identify temporal conflicts, determine whether a planned activity is affected by an anomaly, or reconstruct the sequence of events leading to an operational condition.

Finally, controlled ablation experiments should be introduced to isolate the contribution of individual FUNES components. Variants removing or modifying the Planner, specialized retrieval agents, Critic, iterative refinement, or Synthesizer could be compared under the same benchmark conditions. Such experiments would allow the architectural contribution of each component to be quantified independently and would provide stronger evidence for the architectural claims than the current end-to-end comparison alone.

\section{Operational Validation}

A final direction concerns validation in more realistic operational settings. The current evaluation measures the correctness of generated responses against a predefined operational ground truth. While this is necessary for controlled benchmarking, it does not directly measure whether the generated information improves the work of human operators.

Future studies should therefore involve domain experts in the evaluation process. Expert operators could assess not only factual correctness, but also operational usefulness, relevance, trustworthiness, clarity of uncertainty, and suitability of proposed actions. Human evaluation would provide a complementary perspective to automated scoring and could identify failure modes that are difficult to capture through predefined evaluation criteria.

The effect of FUNES on the decision-making process could also be evaluated directly. Metrics such as time-to-decision, number of required information queries, number of operator interventions, and rate of correct operational decisions could provide evidence on whether the system produces a measurable benefit beyond response-level accuracy. This would shift the evaluation from assessing the quality of individual answers towards assessing the effectiveness of the overall human--AI decision-support process.

Another relevant research direction is the evaluation of explainability and evidence provenance. A future version of FUNES could explicitly expose the evidence supporting each major conclusion, allowing operators to trace an operational assessment back to the corresponding network measurements, tickets, planning records, or knowledge-base entries. Experiments could then investigate whether this form of evidence-grounded explanation improves operator trust, verification speed, and resistance to unsupported model conclusions.

Finally, FUNES should be evaluated under more dynamic operating conditions. The current benchmark is based on a controlled snapshot of operational information, whereas real ground-segment environments involve continuously changing network states, planning updates, newly generated tickets, and evolving investigations. Future experiments could therefore evaluate long-running investigations in which the system maintains state across multiple queries and updates its assessment as new evidence becomes available. This would provide a more realistic validation of the shared state and graph-based orchestration mechanisms and would assess whether FUNES can maintain consistency across successive investigation cycles.

Overall, these research directions would extend the current evaluation from a controlled comparison of two response-generation configurations towards a broader investigation of information orchestration, LLM reasoning, evidence-grounded decision support, and human--AI interaction. Such extensions would make it possible to determine not only whether FUNES improves response quality in the considered benchmark, but also which architectural mechanisms produce the improvement, under which task conditions they remain effective, and whether the resulting benefits translate into measurable improvements in operational decision support.